\begin{ThreePartTable}

\begin{TableNotes}
    \item[1] In this table, LP GPIO No. and GPIO No. are used, not the pin No.
    \item[2] GPIO12 and GPIO13 are not LP GPIO.
\end{TableNotes}
   \begin{longtable}[c]{|p{2cm}|p{2cm}|p{2.5cm}|p{3cm}|p{3cm}| }
   \caption{Analog Functions of IO MUX Pins}
    \label{table:rtcmux_analog_pads} \\
        \hline
        \rowcolor{lightgray}
        \textbf{LP GPIO No.\tnote{1}} &\textbf{GPIO No.\tnote{1}} & \textbf{Pin Name} & \textbf{Analog Function 0} & \textbf{Analog Function 1} \\
    \hline 
    \endfirsthead
    \insertTableNotes
    \endlastfoot
 0 & 0 & XTAL\_32K\_P   & XTAL\_32K\_P & ADC1\_CH0                  \\
\hline
 1 & 1 & XTAL\_32K\_N   & XTAL\_32K\_N & ADC1\_CH1                 \\
\hline
 2 & 2 & GPIO2   & -            & ADC1\_CH2                 \\
\hline
 3 & 3 & GPIO3   & -            & ADC1\_CH3                 \\
\hline
 4 & 4 & MTMS   & -            & ADC1\_CH4                  \\
\hline
 5 & 5 & MTDI   & -            & ADC1\_CH5                 \\
\hline
 6 & 6 & MTCK   & -            & ADC1\_CH6                 \\
\hline
 - & 12 & GPIO12\tnote{2}   & USB\_D-            & -                 \\
\hline
 - & 13 & GPIO13\tnote{2}   & USB\_D+            & -                 \\
\hline


\end{longtable}

\end{ThreePartTable}

